\chapter{Ciguatera poisoning}
\index{Ciguatera poisoning}

\section*{Definition and Overview}
Ciguatera poisoning (also known as ciguatera fish poisoning or CFP) is a distinctive, globally prevalent foodborne illness caused by the consumption of warm-water reef fish containing ciguatoxins. These potent neurotoxins are produced by benthic dinoflagellates of the genus \textit{Gambierdiscus}, which proliferate on coral reefs. The toxins enter the marine food web when herbivorous fish graze on the algae, and they subsequently undergo biomagnification as they pass up the trophic levels to larger predatory fish. Ciguatoxins are heat-stable, lipid-soluble, acid-resistant, and unaffected by cooking, freezing, or gastric acid. Clinical presentation is characterised by a unique combination of acute gastrointestinal, neurological, and cardiovascular symptoms, notably including cold allodynia (a sensory reversal where cold stimuli are perceived as burning hot).

\section*{Epidemiology}
Ciguatera poisoning is the most common non-bacterial foodborne illness related to fish consumption worldwide, affecting an estimated 20,000 to 50,000 people annually, though under-reporting is widespread. It is endemic to tropical and subtropical regions, particularly the Pacific Ocean, the Indian Ocean, and the Caribbean Sea. In many countries, ciguatera poisoning is most frequently associated with fish caught in the warm coastal waters of Queensland and the Northern Territory, such as Spanish mackerel, coral trout, red bass, and chinamanfish. Climate change and rising sea temperatures are expanding the geographical range of \textit{Gambierdiscus} species, leading to emerging cases in historically cooler temperate zones. International trade in tropical seafood also contributes to cases in non-endemic areas.

\section*{Aetiology and Pathophysiology}
The aetiology of ciguatera poisoning is directly linked to the ingestion of ciguatoxins, which exert profound effects on excitable membranes in the human body.
\begin{itemize}
    \item \textbf{Toxin production and bioaccumulation:} Benthic dinoflagellates produce precursor toxins (gambiertoxins) that are biotransformed by fish into highly toxic ciguatoxins (CTXs), which accumulate in fish tissues, particularly the liver, viscera, and head.
    \item \textbf{Sodium channel activation:} Ciguatoxins bind selectively to voltage-gated sodium channels (specifically site 5) on nerve and muscle membranes, keeping them open at resting potentials. This leads to continuous sodium influx, repetitive membrane depolarisation, and cell swelling.
    \item \textbf{Neurological hyperexcitability:} Persistent sodium channel activation causes spontaneous firing of peripheral nerves, leading to sensory abnormalities such as paraesthesia and dysaesthesia.
    \item \textbf{Autonomic dysfunction:} The toxin disrupts parasympathetic and sympathetic pathways, which can result in cardiovascular instability, including profound bradycardia and arterial hypotension.
    \item \textbf{Immunological and inflammatory response:} Ciguatoxins may also trigger histamine release from mast cells, contributing to the severe pruritus (itching) commonly reported by patients.
\end{itemize}

\section*{Clinical Features}
Symptoms typically develop within 2 to 24 hours of fish ingestion, presenting in distinct gastrointestinal, cardiovascular, and neurological phases.
\begin{itemize}
    \item \textbf{Gastrointestinal symptoms:} Early manifestations include watery diarrhoea, abdominal cramps, nausea, vomiting, and dehydration, typically resolving within 24 to 48 hours.
    \item \textbf{Cold allodynia:} A hallmark neurological symptom where cold objects or liquids produce a painful, burning, or electric shock sensation, while hot objects may feel cold.
    \item \textbf{Neurological disturbances:} Widespread paraesthesia (tingling or numbness) of the lips, tongue, perioral area, hands, and feet, accompanied by circumoral pruritus, severe headache, vertigo, and profound generalised fatigue.
    \item \textbf{Cardiovascular signs:} Bradycardia (slow heart rate), hypotension (low blood pressure), and orthostatic intolerance, which can occur in severe cases due to autonomic disruption.
    \item \textbf{Musculoskeletal pain:} Myalgia, arthralgia, and muscle weakness, particularly in the lower limbs, making ambulation difficult.
    \item \textbf{Chronic neurological sequelae:} In a subset of patients, neurological symptoms (such as fatigue, paresthesias, and itching) can persist or recur for months or even years, often exacerbated by alcohol, nuts, or caffeine consumption.
\end{itemize}

\section*{Differential Diagnosis}
Ciguatera poisoning must be distinguished from other marine toxin syndromes and systemic neurological conditions.
\begin{itemize}
    \item \textbf{Scombroid poisoning:} Caused by histamine buildup in improperly refrigerated dark-meat fish (e.g., tuna), presenting with rapid-onset facial flushing, hives, and throbbing headache, but lacking the neurological symptoms of ciguatera.
    \item \textbf{Paralytic shellfish poisoning:} Caused by saxitoxin in contaminated shellfish, presenting with rapid, progressive flaccid paralysis and respiratory depression.
    \item \textbf{Neurotoxic shellfish poisoning:} Caused by brevetoxins, which also target sodium channels and produce temperature reversal, but are typically associated with shellfish consumption and aerosol exposure during red tides.
    \item \textbf{Guillain-Barr\'{e} syndrome:} An autoimmune demyelinating polyneuropathy presenting with ascending muscle weakness and paraesthesia, but without the acute gastrointestinal onset or link to fish consumption.
    \item \textbf{Organophosphate poisoning:} Ingestion of pesticide-contaminated food, causing cholinergic excess (salivation, lacrimation, urination, diarrhoea, bradycardia), which can be differentiated by occupational history or specific laboratory testing.
\end{itemize}

\section*{When to Seek Medical Care}
Patients who develop neurological symptoms or severe vomiting and diarrhoea after eating tropical or subtropical fish should seek prompt medical evaluation.

\textbf{Contact your local emergency services for:}
\begin{itemize}
    \item Severe difficulty breathing, chest pain, or a feeling of throat tightness.
    \item Severe dizziness, confusion, inability to stand, or loss of consciousness.
    \item Extreme bradycardia or a dangerously low blood pressure.
    \item Progressive muscle weakness or paralysis.
\end{itemize}

\section*{Diagnosis and Investigations}
Diagnosis is primarily clinical, based on the classic presentation of symptoms and a history of reef fish ingestion.
\begin{itemize}
    \item \textbf{Clinical history:} Documenting the type of fish consumed, geographic origin of the catch, time of onset, and presence of temperature reversal.
    \item \textbf{Serum electrolytes and renal function:} Essential to monitor and manage dehydration and electrolyte imbalances resulting from severe diarrhoea and vomiting.
    \item \textbf{Electrocardiogram:} Recommended to monitor for bradycardia, conduction blocks, or arrhythmias in patients demonstrating cardiovascular symptoms.
    \item \textbf{Toxin detection in fish remnants:} If leftover fish is available, specialized reference laboratories can perform liquid chromatography-tandem mass spectrometry (LC-MS/MS) or cell-based assays to detect ciguatoxins. There is currently no widely available diagnostic test to detect ciguatoxins in human blood or tissue.
\end{itemize}

\section*{Management}
There is no specific antidote for ciguatera poisoning; management is primarily supportive and symptomatic.
\begin{itemize}
    \item \textbf{Intravenous fluid resuscitation:} Restoring intravascular volume and correcting electrolyte imbalances in patients with severe gastrointestinal losses.
    \item \textbf{Intravenous mannitol:} Historically administered (e.g., 1\,g/kg of 20\% mannitol infused over 30--45 minutes) within the first 48--72 hours of symptom onset to alleviate severe neurological symptoms, although clinical trials show mixed efficacy.
    \item \textbf{Analgesics and neuropathic agents:} Gabapentin, amitriptyline, or non-steroidal anti-inflammatory drugs (NSAIDs) may be prescribed to manage persistent neuropathic pain and paraesthesias.
    \item \textbf{Antihistamines:} Medications like hydroxyzine or cetirizine to manage severe pruritus.
    \item \textbf{Dietary restrictions:} Patients must strictly avoid alcohol, fish, shellfish, nuts, and caffeine for several weeks to months following recovery, as these substances can trigger a recurrence of neurological symptoms.
\end{itemize}

\section*{Complications}
\begin{itemize}
    \item \textbf{Severe dehydration and shock:} Resulting from uncompensated gastrointestinal fluid loss.
    \item \textbf{Cardiovascular collapse:} Secondary to severe bradycardia and hypotension in highly toxic ingestions.
    \item \textbf{Chronic ciguatera syndrome:} Persistent, debilitating neurological symptoms, fatigue, and depression lasting for months, impacting functional capacity.
    \item \textbf{Aspiration pneumonia:} Risk in patients with impaired airway protection due to severe neurological impairment or vomiting.
\end{itemize}

\section*{Prognosis and Outcomes}
Most patients with ciguatera poisoning recover fully within a few days to several weeks. However, the prognosis is highly variable, and neurological recovery can be protracted. Up to 20--30\% of affected individuals experience chronic symptoms that recur intermittently, often triggered by dietary triggers (particularly alcohol or seafood). The mortality rate is low (typically $<1$\%), with rare fatalities occurring due to respiratory failure from respiratory muscle paralysis or cardiovascular collapse.

\section*{Prevention and Self-Care}
\begin{itemize}
    \item \textbf{Avoid high-risk fish species:} Limit or avoid eating large, predatory reef fish (such as barracuda, moray eels, large mackerel, and red bass) caught in endemic areas.
    \item \textbf{Inquire about fish origin:} When purchasing fish or dining in tropical areas, ask where the fish was caught and avoid fish from known hot spots.
    \item \textbf{Do not eat fish viscera:} Avoid consuming the liver, head, roe, and intestines of reef fish, as ciguatoxins are highly concentrated in these organs.
    \item \textbf{Avoid test-and-eat assumptions:} Do not rely on local myths (such as feeding fish to ants or domestic pets, or checking if silver tarnishes) to determine if a fish is safe.
    \item \textbf{Avoid dietary triggers post-exposure:} If previously poisoned, avoid alcohol, nuts, and fish to prevent triggering symptom recurrence.
\end{itemize}

\section*{Sources and Further Reading}
The following reputable organisations provide further information on ciguatera poisoning:
\begin{itemize}
    \item Queensland Health. \textit{Fact sheets and safety guidelines on ciguatera fish poisoning.} https://www.health.qld.gov.au/
    \item Healthdirect Australia. \textit{Overview of ciguatera food poisoning symptoms and prevention.} https://www.healthdirect.gov.au/
    \item World Health Organization (WHO). \textit{Information on natural toxins in food and seafood safety.} https://www.who.int/
    \item Centers for Disease Control and Prevention (CDC). \textit{Traveler's health guide to ciguatera fish poisoning.} https://www.cdc.gov/
    \item Australian Institute of Marine Science. \textit{Research on Gambierdiscus distribution and climate impacts.} https://www.aims.gov.au/
\end{itemize}
